%%%%%%%%%%%%%%%%%
% This is an sample CV template created using altacv.cls
% (v1.7, 9 August 2023) written by LianTze Lim (liantze@gmail.com). Compiles with pdfLaTeX, XeLaTeX and LuaLaTeX.
%
%% It may be distributed and/or modified under the
%% conditions of the LaTeX Project Public License, either version 1.3
%% of this license or (at your option) any later version.
%% The latest version of this license is in
%%    http://www.latex-project.org/lppl.txt
%% and version 1.3 or later is part of all distributions of LaTeX
%% version 2003/12/01 or later.
%%%%%%%%%%%%%%%%

%% Use the "normalphoto" option if you want a normal photo instead of cropped to a circle
% \documentclass[10pt,a4paper,normalphoto]{altacv}

\documentclass[9pt,a4paper,ragged2e,withhyper]{altacv}
%% AltaCV uses the fontawesome5 and packages.
%% See http://texdoc.net/pkg/fontawesome5 for full list of symbols.

% Change the page layout if you need to
\geometry{left=0.8cm,right=0.8cm,top=1.1cm,bottom=1.1cm,columnsep=1.2cm}

% The paracol package lets you typeset columns of text in parallel
\usepackage{paracol}

\usepackage{twemojis}
% Change the font if you want to, depending on whether
% you're using pdflatex or xelatex/lualatex
% WHEN COMPILING WITH XELATEX PLEASE USE
% xelatex -shell-escape -output-driver="xdvipdfmx -z 0" sample.tex
\ifxetexorluatex
  % If using xelatex or lualatex:
  \setmainfont{Roboto Slab}
  \setsansfont{Lato}
  \renewcommand{\familydefault}{\sfdefault}
\else
  % If using pdflatex:
  \usepackage[rm]{roboto}
  \usepackage[defaultsans]{lato}
  % \usepackage{sourcesanspro}
  \renewcommand{\familydefault}{\sfdefault}
\fi

% Change the colours if you want to
\definecolor{SlateGrey}{HTML}{2E2E2E}
\definecolor{LightGrey}{HTML}{666666}
\definecolor{DarkPastelRed}{HTML}{450808}
\definecolor{PastelRed}{HTML}{8F0D0D}
\definecolor{GoldenEarth}{HTML}{E7D192}
\colorlet{name}{black}
\colorlet{tagline}{PastelRed}
\colorlet{heading}{DarkPastelRed}
\colorlet{headingrule}{GoldenEarth}
\colorlet{subheading}{PastelRed}
\colorlet{accent}{PastelRed}
\colorlet{emphasis}{SlateGrey}
\colorlet{body}{LightGrey}

% Change some fonts, if necessary
\renewcommand{\namefont}{\Huge\rmfamily\bfseries}
\renewcommand{\personalinfofont}{\footnotesize}
\renewcommand{\cvsectionfont}{\LARGE\rmfamily\bfseries}
\renewcommand{\cvsubsectionfont}{\large\bfseries}


% Change the bullets for itemize and rating marker
% for \cvskill if you want to
\renewcommand{\cvItemMarker}{{\small\textbullet}}
\renewcommand{\cvRatingMarker}{\faCircle}
% ...and the markers for the date/location for \cvevent
% \renewcommand{\cvDateMarker}{\faCalendar*[regular]}
% \renewcommand{\cvLocationMarker}{\faMapMarker*}


% If your CV/résumé is in a language other than English,
% then you probably want to change these so that when you
% copy-paste from the PDF or run pdftotext, the location
% and date marker icons for \cvevent will paste as correct
% translations. For example Spanish:
% \renewcommand{\locationname}{Ubicación}
% \renewcommand{\datename}{Fecha}


\begin{document}
\name{Nathan Cassereau}
\tagline{AI DevTech @ NVIDIA}
%% You can add multiple photos on the left or right
\photoR{2.8cm}{moi}
% \photoL{2.5cm}{Yacht_High,Suitcase_High}

\NewInfoField{website}{\faGlobe}[https://]
\NewInfoField{age}{\faCalendarDay}[]
\NewInfoField{nationality}{\twemoji{flag: France} }[]

\personalinfo{%
  \email{nathan@cassereau.fr}
  \phone{+33 6 77 54 14 95}
  % \mailaddress{214 Boulevard du Moulin de la Tour, A104, 92140 Clamart (FRANCE)}
  \nationality{French nationality}
  \age{28 years old}
  \location{Clamart, FRANCE}
  \linkedin{nathan-cassereau}
  \github{ncassereau}
  \website{cassereau.fr}
  
  %% You can add your own arbitrary detail with
  %% \printinfo{symbol}{detail}[optional hyperlink prefix]
  % \printinfo{\faPaw}{Hey ho!}[https://example.com/]

  %% Or you can declare your own field with
  %% \NewInfoFiled{fieldname}{symbol}[optional hyperlink prefix] and use it:
  % \NewInfoField{gitlab}{\faGitlab}[https://gitlab.com/]
  % \gitlab{your_id}
  %%
  %% For services and platforms like Mastodon where there isn't a
  %% straightforward relation between the user ID/nickname and the hyperlink,
  %% you can use \printinfo directly e.g.
  % \printinfo{\faMastodon}{@username@instace}[https://instance.url/@username]
  %% But if you absolutely want to create new dedicated info fields for
  %% such platforms, then use \NewInfoField* with a star:
  % \NewInfoField*{mastodon}{\faMastodon}
  %% then you can use \mastodon, with TWO arguments where the 2nd argument is
  %% the full hyperlink.
  % \mastodon{@username@instance}{https://instance.url/@username}
}

\makecvheader
%% Depending on your tastes, you may want to make fonts of itemize environments slightly smaller
% \AtBeginEnvironment{itemize}{\small}

%% Set the left/right column width ratio to 6:4.
\columnratio{0.6}

Senior HPC/AI engineer with a track record in \textbf{GPU kernel optimization} (CUDA, Triton: \textbf{19× speedup, 5× memory reduction}), \textbf{distributed training} at scale, and ML research collaboration on national supercomputers (\textbf{Jean Zay, GB200 NVL72}). Joining \textbf{NVIDIA} in October 2026 as \textbf{AI DevTech}.


% Start a 2-column paracol. Both the left and right columns will automatically
% break across pages if things get too long.
\cvsection{Experience}

\cvevent{AI DevTech}{NVIDIA}{October 2026 (ongoing)}{Courbevoie, FRANCE}
Joining NVIDIA's AI Developer Technologies team, focused on GPU/AI ecosystem enablement for research and industry.

\divider

\cvevent{Senior HPC AI Engineer}{AMIAD}{October 2025 -- October 2026}{Palaiseau, FRANCE}
\begin{itemize}
    \item \textbf{Kernel Optimization:} Developed custom \textbf{CUDA} and \textbf{Triton} kernels achieving significant performance improvements:
    \begin{itemize}
        \item Second-order optimizer: \textbf{19× speedup} and \textbf{5× memory reduction} vs. PyTorch on B200
        \item Custom CUDA kernel: \textbf{4.5× faster} than JAX/XLA compiler
    \end{itemize}
    \item \textbf{HPC Infrastructure Support:} Supporting researchers on \textbf{NVIDIA NVL72} supercomputer (Spack, Slurm scripting, cluster management)
    \item Built \textbf{OhMyJob}, a real-time monitoring platform for Slurm HPC jobs (CPU/IO/GPU metrics collected on compute nodes, centralised into dashboards for researchers)
    \item Infrastructure deployment and automation using \textbf{Docker}, \textbf{Ansible}, and \textbf{OpenStack} for research computing environments.
\end{itemize}
\divider

\cvevent{ML Engineer (Freelance)}{Entalpic}{June 2025 -- September 2025}{Paris, FRANCE}
I optimized their use of \textbf{national supercomputers}, achieving
\textbf{5× throughput improvement} in distributed training. I also implemented specialized
\textbf{MLIP benchmarks} and optimized training pipelines for chemistry AI models.

\divider

\cvevent{Artificial Intelligence Engineer}{IDRIS (CNRS)}{May 2021 -- May 2025}{Orsay, FRANCE}
\begin{itemize}
\item \textbf{AI support team:} supporting \textbf{2000+ researchers} on \textbf{Jean Zay supercomputer}, benchmarking hardware accelerators (NVIDIA/AMD GPUs, GraphCore IPU), developing tools for \textbf{distributed training}
\item \textbf{Research collaborations} (PNRIA): AION (Flatiron Institute, \textbf{foundational model} for astrophysics), AutoSynth (LaBRI, \textbf{LLM code generation})
\item \textbf{Teaching:} Deep Learning courses (Intro, \textbf{Optimization for Supercomputers}, \textbf{LLM Specialization}), NVIDIA bootcamps, JSALT workshops
\end{itemize}

\cvsection{Education}
\setlength{\columnsep}{0cm}
\columnratio{0.55}
\begin{paracol}{2}
\cvevent{MSc.\ Artificial Intelligence \& Deep Learning with Honours}{Imperial College London}{2019 -- 2020}{London, UNITED KINGDOM}
\switchcolumn
\cvevent{Master in Engineering}{CentraleSupélec}{2017 -- 2020}{Gif-sur-Yvette, FRANCE}
\end{paracol}

\columnratio{0.3}
\setlength{\columnsep}{1cm}
\begin{paracol}{2}
\cvsection{Languages}
\cvskill{English}{5}
\divider
\cvskill{French}{5}

\switchcolumn

\cvsection{Skills \& Technologies}

\cvtag{\textbf{GPU \& Kernel Optimization} (CUDA, Triton, GPU Profiling: Nsight / Torch)}

\cvtag{\textbf{Distributed \& Parallel Computing} (OpenMP, PyTorch DDP/FSDP, DeepSpeed, Megatron)}

\cvtag{\textbf{ML Frameworks \& Tools} (PyTorch, JAX, HuggingFace, vLLM, W\&B, MLflow)}

\cvtag{\textbf{Deep Learning} (NLP \& LLM, Computer Vision, Diffusion, Multi-modal, Post-training)}

\cvtag{\textbf{HPC Infrastructure} (Slurm, Spack, Docker, Apptainer, Modules, Ansible, Terraform)}

\cvtag{\textbf{Programming \& Tools} (Python, C++, Rust, Jupyter, Git, CI/CD, \LaTeX, Grafana)}

\end{paracol}

\end{document}
